\section{Marco contextual}
\label{sec:marco_contextual}

La investigación se desarrolla en un contexto de educación media colombiana, con estudiantes de grado décimo y undécimo en una institución educativa de carácter oficial. El interés del estudio se centra en escenarios de aula reales, con condiciones habituales de tiempo, infraestructura y heterogeneidad académica.

En este contexto, la enseñanza de la trigonometría suele concentrarse en unidades curriculares de corta duración y con presión por cobertura de contenidos, lo que frecuentemente reduce el espacio para argumentación y análisis de errores. Esta condición contextual refuerza la necesidad de diseñar tareas que, aun en tiempos limitados, sostengan exigencia matemática y discusión matemática.

También se reconoce una incorporación creciente de recursos digitales en la escuela, aunque con diferencias en acceso, conectividad y familiaridad tecnológica entre estudiantes. Por ello, la integración de IA en esta investigación se plantea de forma gradual, con actividades guiadas y criterios didácticos explícitos, evitando depender de condiciones tecnológicas excepcionales.

El contexto institucional, curricular y tecnológico descrito no se considera un telón de fondo secundario, sino una variable que incide directamente en la viabilidad del diseño de tareas y en la interpretación de los resultados del estudio.
